% GENERATED FILE. Edit canonical lessons, then run npm run generate:books.
% canonical-source-hash: fnv1a64:3edc926ae119c0f6
% canonical-lessons: TA-C87-tiredness, TA-C87-cold, TA-C87-cough, TA-C87-weariness, TA-C87-rest

\chapter{Not Feeling Well}
\label{ch:ta-not-feeling-well}
\hlchaptermodality{87}

\begin{chapteropening}
\textbf{By the end of this chapter:} \emph{I can say I am tired or weary, have a cold or a cough, and tell someone to rest.}

Complete the last lesson of chapter 87: i can say I am tired or weary, have a cold or a cough, and tell someone to rest.
\end{chapteropening}

\section[kaḷaippu]{\ta{களைப்பு} (kaḷaippu) — tiredness}

\label{lesson:TA-C87-tiredness}



\begin{quote}
Before the new one: say the Tamil for hunger, then the Tamil for sleep.
\end{quote}

\subsection*{You'll want to know: \ta{களைப்பு}}
\textbf{\ta{களைப்பு}} (\emph{kaḷaippu}) — \textquotedblleft{}tiredness\textquotedblright{}.

\textbf{\ta{எனக்குக்} \ta{களைப்பாக} \ta{இருக்கிறது}} — \textquotedblleft{}I feel tired\textquotedblright{}, literally \textquotedblleft{}to me it is tired-ly\textquotedblright{}. The person sits in \textbf{\ta{எனக்கு}}, \textquotedblleft{}to me\textquotedblright{}.

\begin{grammarlens}[title={What you've built}]
A first word for how the body feels.
\end{grammarlens}

\subsection*{Your turn}
\begin{itemize}
\raggedright
  \item \emph{Say it:} \emph{kaḷaippu}
  \item \emph{Say it:} \emph{kaḷaippu}, once more
  \item \emph{Say it:} say \emph{eṉakkuk kaḷaippāka irukkiṟatu}
  \item \emph{Recall:} say the Tamil for hunger, then the Tamil for sleep, then say \emph{kaḷaippu} again
\end{itemize}

\begin{tcolorbox}[breakable,colback=teal!4,colframe=teal!35!black,title={Before you move on}]
What does \ta{களைப்பு} mean? (\textquotedblleft{}Tiredness\textquotedblright{}.) Say it once more.
\end{tcolorbox}
\section[caḷi]{\ta{சளி} (caḷi) — a cold, a runny nose}

\label{lesson:TA-C87-cold}



\begin{quote}
Before the new one: say the Tamil for sleep, then the Tamil for tiredness.
\end{quote}

\subsection*{You'll want to know: \ta{சளி}}
\textbf{\ta{சளி}} (\emph{caḷi}) — \textquotedblleft{}a cold\textquotedblright{}, the kind with a runny nose.

\textbf{\ta{எனக்குச்} \ta{சளி} \ta{பிடித்திருக்கிறது}} — \textquotedblleft{}I have caught a cold\textquotedblright{}, literally \textquotedblleft{}a cold has caught me\textquotedblright{}.

\begin{grammarlens}[title={What you've built}]
An illness that catches you.
\end{grammarlens}

\subsection*{Your turn}
\begin{itemize}
\raggedright
  \item \emph{Say it:} \emph{caḷi}
  \item \emph{Say it:} \emph{caḷi}, once more
  \item \emph{Say it:} say \emph{eṉakkuc caḷi piṭittirukkiṟatu}
  \item \emph{Recall:} say the Tamil for sleep, then the Tamil for tiredness, then say \emph{caḷi} again
\end{itemize}

\begin{tcolorbox}[breakable,colback=teal!4,colframe=teal!35!black,title={Before you move on}]
What does \ta{சளி} mean? (\textquotedblleft{}A cold, a runny nose\textquotedblright{}.) Say it once more.
\end{tcolorbox}
\section[irumal]{\ta{இருமல்} (irumal) — a cough}

\label{lesson:TA-C87-cough}



\begin{quote}
Before the new one: say the Tamil for tiredness, then the Tamil for a cold.
\end{quote}

\subsection*{You'll want to know: \ta{இருமல்}}
\textbf{\ta{இருமல்}} (\emph{irumal}) — \textquotedblleft{}a cough\textquotedblright{}.

\textbf{\ta{இருமல்} \ta{இருக்கிறது}} — \textquotedblleft{}there is a cough\textquotedblright{}, the plain way to tell a doctor.

\begin{grammarlens}[title={What you've built}]
A second complaint.
\end{grammarlens}

\subsection*{Your turn}
\begin{itemize}
\raggedright
  \item \emph{Say it:} \emph{irumal}
  \item \emph{Say it:} \emph{irumal}, once more
  \item \emph{Say it:} say \emph{irumal irukkiṟatu}
  \item \emph{Recall:} say the Tamil for tiredness, then the Tamil for a cold, then say \emph{irumal} again
\end{itemize}

\begin{tcolorbox}[breakable,colback=teal!4,colframe=teal!35!black,title={Before you move on}]
What does \ta{இருமல்} mean? (\textquotedblleft{}A cough\textquotedblright{}.) Say it once more.
\end{tcolorbox}
\section[cōrvu]{\ta{சோர்வு} (cōrvu) — weariness, fatigue}

\label{lesson:TA-C87-weariness}



\begin{quote}
Before the new one: say the Tamil for a cold, then the Tamil for a cough.
\end{quote}

\subsection*{You'll want to know: \ta{சோர்வு}}
\textbf{\ta{சோர்வு}} (\emph{cōrvu}) — \textquotedblleft{}weariness\textquotedblright{}, the drained feeling after an illness.

\textbf{\ta{களைப்பு}} is the tiredness of a long day; \textbf{\ta{சோர்வு}} lasts longer. \textbf{\ta{சோர்வாக} \ta{இருக்கிறது}} — \textquotedblleft{}I feel worn out\textquotedblright{}.

\begin{grammarlens}[title={What you've built}]
Two kinds of tired.
\end{grammarlens}

\subsection*{Your turn}
\begin{itemize}
\raggedright
  \item \emph{Say it:} \emph{cōrvu}
  \item \emph{Say it:} \emph{cōrvu}, once more
  \item \emph{Say it:} say \emph{cōrvāka irukkiṟatu}
  \item \emph{Recall:} say the Tamil for a cold, then the Tamil for a cough, then say \emph{cōrvu} again
\end{itemize}

\begin{tcolorbox}[breakable,colback=teal!4,colframe=teal!35!black,title={Before you move on}]
What does \ta{சோர்வு} mean? (\textquotedblleft{}Weariness, fatigue\textquotedblright{}.) Say it once more.
\end{tcolorbox}
\section[ōyvu]{\ta{ஓய்வு} (ōyvu) — rest}

\label{lesson:TA-C87-rest}



\begin{quote}
Before the new one: say the Tamil for a cough, then the Tamil for weariness.
\end{quote}

\subsection*{You'll want to know: \ta{ஓய்வு}}
\textbf{\ta{ஓய்வு}} (\emph{ōyvu}) — \textquotedblleft{}rest\textquotedblright{}.

\textbf{\ta{ஓய்வு} \ta{எடுங்கள்}} — \textquotedblleft{}please take some rest\textquotedblright{}, what you say to someone who is unwell.

\begin{grammarlens}[title={What you've built}]
Tired, a cold, a cough, weary, and rest.
\end{grammarlens}

\subsection*{Your turn}
\begin{itemize}
\raggedright
  \item \emph{Say it:} \emph{ōyvu}
  \item \emph{Say it:} \emph{ōyvu}, once more
  \item \emph{Say it:} say \emph{ōyvu eṭuṅkaḷ}
  \item \emph{Recall:} say the Tamil for a cough, then the Tamil for weariness, then say \emph{ōyvu} again
\end{itemize}

\begin{tcolorbox}[breakable,colback=teal!4,colframe=teal!35!black,title={Before you move on}]
What does \ta{ஓய்வு} mean? (\textquotedblleft{}Rest\textquotedblright{}.) Say it once more.
\end{tcolorbox}
